
@@M0@@. Soient @@M1@@ les composantes connexes de @@M2@@~; alors @@M3@@ se décompose, @@M4@@, et l'assertion en résulte. Enfin, une récurrence immédiate nous ramène au cas d'un morphisme @@M5@@. @@M6@@ et @@M7@@ résultent alors, quitte à restreindre @@M8@@ d'après (EGA III 7.8.6), de la

\textbf{Proposition 2.3.} \uline{Soient} @@M9@@ \uline{un schéma,} @@M10@@ \uline{un morphisme propre et} @@M11@@ \uline{deux sous-schémas fermés définis par les idéaux} @@M12@@\uline{. On suppose que}
\begin{enumerate}[label=\alph*)]
\item @@M13@@ \uline{sont fidèlement plats et de présentation finie sur} @@M14@@.
\item @@M15@@ et @@M16@@ \uline{universellement (où} @@M17@@, @@M18@@).
\item @@M19@@.
\end{enumerate}
\uline{Alors} @@M20@@ \uline{est représentable par un morphisme affine de présentation finie.}

En effet notons d'abord que les hypothèses sont stables pour tout changement de base. Pour (c), cela résulte des deux suites exactes~:
@@M21@@
et
@@M22@@
la première montre que @@M23@@ est plat sur @@M24@@~; ensuite la seconde que @@M25@@ est stable par changement de base.

L'assertion est locale pour la topologie f.p.p.f. En faisant le changement de base @@M26@@, on peut donc supposer que @@M27@@ possède une section @@M28@@, qui définit des sections @@M29@@ de @@M30@@ et @@M31@@ de @@M32@@.

Après un changement de base convenable, on est ramené à voir que tout point @@M33@@ a comme image réciproque un sous-foncteur de
